% A section suitable for inclusion in the main report.
\section{Scope and substitution audit of the untagged representation}

The concrete representation uses two disjoint classes of variables.
The integers $-1,\ldots,-m$ are free numerical parameters at scope $m$.
The positive integers $1,\ldots,d$ are de Bruijn indices for $d$ locally
bound numerical variables. Zero is the numerical constant, not a variable.
At positive depth $d$, the syntax is typed as follows:
\[
\begin{array}{c|c|c}
\text{object}&\text{representation}&\text{depths of its children}\\\hline
R(n,a;u.h)&[n,a,h]&(d,d,d+1)\\
s=t&[s,t]&(d,d)\\
D_{x,y}B&[B]&(d+2)
\end{array}
\]
An expression may use free parameters only in $[-m,-1]$. In a paired
matrix, $y$ is the nearest variable (index 1) and $x$ is next (index 2).
A canonical recursor whose counter is literally zero must have body zero.
Arrays have no other permitted arities. Exact integer tests exclude
booleans; strings, nulls, floating-point numbers, and arbitrary objects
are not terms.

\subsection{The opening operation and why lifting is unnecessary}

Let $O_d(e,v)$ remove the index $d+1$ while preserving the $d$ nearer
indices. On an integer, it replaces $d+1$ by $v$, decrements an integer
greater than $d+1$, and leaves every other integer unchanged. On an
array it recursively processes its children at the depths in the table.
At each resulting recursor node it erases the body if the resulting
counter is zero.

The replacement $v$ must be a numerical term well scoped at positive
depth zero. It can contain internally bound variables inside its own
recursors, but has no positive variable free relative to its root.
Therefore inserting $v$ under any surrounding binders needs no lifting:
its positive indices already refer exclusively to its own binders, and
its negative parameters are names independent of binder depth. This
is the essential precondition. Allowing a positive free variable in
$v$ without lifting would be incorrect.

\begin{lemma}[Scope preservation]
If $e$ has free scope $m$ and ambient positive depth $k>d$, and $v$ is a numerical term
at positive depth zero and free scope $m$, then $O_d(e,v)$ has the same
sort as $e$, ambient positive depth $k-1$, and free scope $m$.
The operation introduces no new free parameter beyond those in $e,v$.
It produces canonical syntax when its inputs are canonical.
\end{lemma}
\begin{proof}
Structural induction on $e$. The integer cases respectively remove the
selected index, lower an outer index, or preserve a nearer index or
parameter. A locally closed replacement remains well scoped under any
additional ambient binders. The three array cases use exactly their
displayed binder increments. Erasing a recursor body only removes
occurrences; it cannot introduce a variable or violate a scope bound.
\end{proof}

\subsection{Coherence of canonicalization}

Let $\nu$ recursively canonicalize raw expressions by replacing
$R(0,a;u.h)$ by $R(0,a;u.0)$ after normalizing the children. Let
$O_d^{\rm raw}$ be the same opening without body erasure. For canonical
inputs $e,v$, the executable opening satisfies the first identity below;
the second is the normalization-coherence identity on arbitrary raw inputs:
\[
 O_d(e,v)=\nu\bigl(O_d^{\rm raw}(e,v)\bigr),
 \qquad
 \nu\bigl(O_d^{\rm raw}(\nu e,\nu v)\bigr)
 =\nu\bigl(O_d^{\rm raw}(e,v)\bigr).
\]
Both identities follow by structural induction. A body discarded before
opening had a counter already literally zero, which remains zero under
opening. A counter that becomes zero only after opening is discarded on
both sides by the outer normalization. Also $\nu(t)=0$ if and only if
$t$ is literally zero: canonicalization does not collapse an entire
recursor node to a number. Thus normalizing early cannot conceal a
nonzero counter that subsequently becomes a different constructor.

For a matrix with slots $y=1,x=2$ and locally closed terms $t,w$,
\[
 O_0(O_1(B,t),w)
 =\nu\bigl(B[t/x,w/y]\bigr).
\]
The two openings may also be performed in the opposite order, opening
$y$ first and then the newly nearest $x$. Their equality follows from
the ordinary simultaneous-substitution lemma, since neither replacement
contains a free positive index, followed by the displayed normalization
coherence. This identifies the executable operation with the mathematical
normalized substitution, including nested recursors and nested $D$.

\subsection{Each use in the five-rule verifier}

At free scope $m$, rule inputs and generated objects have these types:
\begin{itemize}
\item A public goal is a formula at positive depth zero. A computation
counter and initializer are terms at depth zero, while its body is a
term at depth one. Its right side is $O_0(h,R(n,a;u.h))$.
\item The single transport template is a term at depth one. Both
equation endpoints supplied as images are terms at depth zero, so
the two template instances are valid openings.
\item A context clause is a pair $[A;y.B]$, with $A$ a formula at depth
zero and $B$ a formula at depth one. Clause use supplies a locally
closed target term $t$ and concludes $O_0(B,t)$ after proving $A$.
\item An ordinary assumption $F$ becomes $[H;v.F]$. Since $F$ has no
positive index free at its root, the identical syntax can be viewed
under the unused additional binder $v$; no shift is needed.
\item For paired elimination, use the fresh free parameter
$b=-(m+1)$. From a paired matrix $B$ at depth two, form
$C=O_1(B,b)$ at depth one and retain the clause $[O_0(C,0);y.C]$.
The branch runs at scope $m+1$.
\item For paired introduction, validate the witness $t$ at the original
scope $m$, form $C=O_1(B,t)$, and use the fresh parameter
$z=-(m+1)$. Its ordinary branch assumption and target are respectively
$O_0(C,\eta)$ and $O_0(C,[0,z,0])$, where $\eta=z$ in the progress
mode and zero otherwise. The array $[0,z,0]$ is the canonical term
$Sz$.
\end{itemize}

Old contexts need no shifting when the free scope increases: their
negative integers retain the same parameter meanings. The newly chosen
negative integer is absent from every previously validated input. A
branch can use it, but its discharged result cannot leak it: the outer
goal and the introduction witness are validated at the old scope, and
the branch conclusion must equal the required outer-scope conclusion
or the explicitly constructed target. Independent branches may reuse
the same next negative integer because their contexts are separate.

\paragraph{Verifier boundary.}
Externally supplied goals, witnesses, templates, and computation data
must be validated at their specified scopes before any opening. Generated
clauses inherit their validity from the preceding rules and the scope
lemma. In particular, a positive integer must never be accepted as a
free numerical witness merely because it is an integer, and an
eigenparameter may not appear in an outer goal. A validator which first
erases arbitrary malformed zero-counter bodies may define a different
accepted encoding; strict canonical validation must reject such bodies
before interpreting them.

The arithmetic normalization theorem used here is exactly the inherited
canonical-recursor theorem. This representation adds no new computation
law, equality principle, induction instance, or numerical object.

\paragraph{Clause pairs are a separate data structure.}
A retained pair $(A,B)$ has two components like an equality array, but its
second component has an extra binder. It is never passed to the generic
expression opener as one object. Clause use opens $B$ alone. This
separation is necessary for the stated binding invariant.
